\section{Experimental framework and reading conventions}

\subsection{Reference configuration}

The later experiments use vectorial NA-0.8 excitation, linear fluorescence, a homogeneous medium and negligible specimen-induced scattering. Excitation and detection wavelengths are approximated as equal at 532 nm. A phase plate is propagated to an intermediate illumination plane and relayed into the specimen; the detection focus remains fixed while input angle changes.

Optical propagation uses a 384\ensuremath{\times}384 grid at 0.08 \ensuremath{\mu}m. The analysis object is 32\ensuremath{\times}96\ensuremath{\times}96 at 0.32 \ensuremath{\mu}m, and reconstruction uses 16\ensuremath{\times}48\ensuremath{\times}48 voxels at 0.64 \ensuremath{\mu}m. The represented volume is 30.72\ensuremath{\times}30.72\ensuremath{\times}10.24 \ensuremath{\mu}m. The original detector has NA 0.1 and 12\ensuremath{\times}12 samples per image. A single phase plate has 0.32 \ensuremath{\mu}m phase features and a 6.6 \ensuremath{\mu}m final gap. These are simulation and reconstruction sampling choices, not measured optical resolution.

The baseline reconstruction uses PnP-ADMM with the generic pretrained DnCNN-17M in three orthogonal slice orientations and a sparsity branch. The extended learned solver replaces the slice denoiser and adds a 3D denoising branch. Both use a 20-iteration reconstruction budget. Baseline bead tests in the original package used sparsity alone; the later learned-prior bead control instead compares PnP solvers using their neuronal-validation settings.

\subsection{Measurements, noise and fair comparisons}

Let \(P\) be detector samples per image, \(k\) illumination directions, \(M=Pk\) scalar measurements and \(N=36{,}864\) reconstructed values. Measurement compression is \(C=N/M\). At the original 12\ensuremath{\times}12 detector, \(k=256,64,32,16,8\) corresponds to \(C=1,4,8,16,32\).

The final comparisons use the relative-Gaussian convention
\[
q(M)=15+5\log_{10}(N/M),\qquad
\sigma(M)=\operatorname{mean}(\mathbf y_N^0)\,10^{-q(M)/10}.
\]
Here \(\mathbf y_N^0\) is the relevant encoder/detector's noise-free reference stack at M=N. At a fixed detector this gives the square-root count scaling used in the main work. FRIS calls it the total-dose or equal-detected-photon convention. These names identify a normalisation assumption, not calibrated photon counts or a Poisson detector model. Earlier fixed-source-power and equal-illumination studies are labelled separately because their normalisation changes device rankings.

Image-count reduction and measurement compression must not be confused. A 48\ensuremath{\times}48 detector with 16 images and a 12\ensuremath{\times}12 detector with 256 images both supply 36,864 scalar measurements. The former uses sixteen times fewer images, but both have \(C=1\).

\subsection{Test sets, tuning and evidence levels}

The final quantitative test set contains eight prespecified mouse pyramidal neuron subvolumes with dendritic arbors: p64, p128, p192, p320, p448, p576, p704 and p832. Repeated noise realisations and illumination subsets are averaged within each subvolume before computing across-subvolume statistics. Subvolume 0 and the reference-selected qualitative examples p384/p973 are reported separately. PSNR, foreground SSIM and local depth error measure different aspects of recovery; depth error is a centroid-based metric, not axial resolution.

Extended solver and detector comparisons are tuned on held-out neuronal validation subvolumes p303, p10 and p693. This differs from the main manuscript's earlier vascular tuning. Each solver/encoder/detector/measurement setting is tuned independently, with boundary checks on the search grid. Consequently, extended baseline rows are not automatically interchangeable with the manuscript's baseline rows.

Evidence has three levels in this record. Final results are drawn from the consolidated 70-condition master table and paired comparisons. Earlier reconstruction results retain their original configuration and object labels. Historical surrogate scores describe what guided exploration; they are not final PnP results or measured resolution. The source index at the end identifies the relevant logs, scripts and exports.
